% ======================================================================
%  Supplemental Material
%  Nucleation density programs the viscoelastic gel point of a percolation network
%  REVTeX 4.2
% ======================================================================
\documentclass[aps,prl,onecolumn,superscriptaddress,nofootinbib]{revtex4-2}

\usepackage{graphicx}
\usepackage{amsmath,amssymb}
\usepackage{bm}
\usepackage{booktabs}
\usepackage[colorlinks=true,allcolors=blue]{hyperref}

\graphicspath{{figures_v2/}{./}}

% S-prefixed numbering for SM
\renewcommand{\thesection}{S\Roman{section}}
\renewcommand{\thefigure}{S\arabic{figure}}
\renewcommand{\thetable}{S\arabic{table}}
\renewcommand{\theequation}{S\arabic{equation}}

\newcommand{\msd}{\langle r^2(t)\rangle}
\newcommand{\pc}{p_c}
\newcommand{\pcp}{p_c'}
\newcommand{\kap}{\kappa}
\newcommand{\nueff}{\nu_{\mathrm{eff}}}

\begin{document}

\title{Supplemental Material:\\ Nucleation density programs the viscoelastic gel point of a percolation network}
\author{A.~Author}
\author{B.~Author}
\author{C.~Author}
\author{D.~Author}
\affiliation{[Affiliation], [Institution], [City], [Country]}
\date{\today}
\maketitle

This Supplemental Material gives the lattice and random-walk model, the nucleation-density
growth algorithm, the finite-size-scaling protocol and its full results, the fixed-$L$
gel-point curve, the frequency-domain (GSER) analysis, robustness checks on the growth rule,
and computational details. Equation, figure, and table numbers are prefixed ``S''.

% ======================================================================
\section{Lattice model and the nucleation-density ($\kap$) growth algorithm}
% ======================================================================
We use a simple cubic lattice of side $L$ with periodic boundary conditions. Each site is
\emph{occupied} (network phase) or \emph{unoccupied} (void/fluid phase); the occupation
probability $p$ sets the target occupied count $N_{\mathrm{target}}=\lfloor pL^3\rceil$.
Lattices are built cumulatively over the $p$-sweep, so the occupied set at higher $p$ contains
that at lower $p$ as a subset---consistent with progressive gelation and ensuring continuity of
$\alpha(p)$.

Growth is governed by a nucleation-density parameter $\kap\in[0,1]$. At each increment, of the
$N_{\mathrm{new}}$ sites to be added,
\begin{equation}
  N_{\mathrm{seeds}}=\big\lfloor(1-\kap)\,N_{\mathrm{new}}\big\rceil
\end{equation}
are placed at uniformly random unoccupied positions (independent nuclei), and the remaining
$N_{\mathrm{new}}-N_{\mathrm{seeds}}$ are grown by 6-neighbor (face-sharing) adjacency from the
combined frontier of all occupied sites. Connectivity is guaranteed by construction: every
occupied site is reachable from a seed through adjacency steps. The limits are transparent:
$\kap=0$ places every new site as an independent random seed (Bernoulli site percolation);
$\kap=1$ plants a single seed and grows a compact connected cluster (Eden growth~\cite{Eden1961}).
Intermediate $\kap$ interpolates between a spatially uniform obstacle field and a clustered,
compact network. Physically $\kap$ is the number of independent nucleation centers per unit
increase in network density.

% ======================================================================
\section{Random-walk engine and observables}
% ======================================================================
Random walkers are initialized on unoccupied sites, modeling microrheology probe particles in
the fluid phase. Each step attempts a move to one of the six nearest neighbors chosen uniformly;
if the target is unoccupied the move is accepted, otherwise the walker dwells at its current site
for a binding time $\mathrm{WT}=\lceil|\mathcal{N}(\mu_{\mathrm{WT}},\sigma_{\mathrm{WT}})|\rceil$
steps, with $\mu_{\mathrm{WT}}=20$, $\sigma_{\mathrm{WT}}=5$. This models transient probe--network
binding and suppresses artificial super-diffusion from immediate reflection. Periodic boundaries
are applied through $x_n=\mathrm{mod}(x_n-1,L)+1$; the unwrapped coordinate is used for the MSD,
\begin{equation}
  \msd=\frac{1}{N_W}\sum_{i=1}^{N_W}\big[\Delta x_i(t)^2+\Delta y_i(t)^2+\Delta z_i(t)^2\big],
\end{equation}
accumulated online by parallel reduction to avoid storing $O(L_W N_W)$ trajectories
($\sim$36~GB at production scale). Production runs use $N_W=3000$ walkers and $N_s=3$ replicates
per condition. The anomalous exponent $\alpha$ is the log--log slope of $\msd$ over the
intermediate window $[L_W/100,\,L_W/10]$, and the gel point is located by linear interpolation of
$\alpha(p)$ through $0.5$,
\begin{equation}
  \pcp=p_i+\frac{0.5-\alpha(p_i)}{\alpha(p_{i+1})-\alpha(p_i)}(p_{i+1}-p_i),
  \qquad \alpha(p_i)>0.5>\alpha(p_{i+1}).
\end{equation}

% ======================================================================
\section{Winter--Chambon criterion and the GSER}
% ======================================================================
At the gel point $\tan\delta=G''/G'=1$, i.e.\ $\delta=45^{\circ}$~\cite{Winter1987,Chambon1987}.
For a power-law MSD $\msd\propto t^{\alpha}$ the GSER~\cite{Mason1995,Mason2000}
\begin{equation}
  G^*(\omega)\approx\frac{k_BT}{\pi a\,\langle r^2(1/\omega)\rangle\,\Gamma(1+\alpha(\omega))}
\end{equation}
gives moduli $G',G''\propto\omega^{\alpha}$ and a frequency-independent phase angle
\begin{equation}
  \delta=\tfrac{\pi}{2}\alpha \quad\Longleftrightarrow\quad \delta(^{\circ})=90\,\alpha,
\end{equation}
so $\delta=45^{\circ}$ is exactly $\alpha=0.5$. Dimensional spectra use $T=293$~K,
$\eta=1.2\times10^{-3}$~Pa\,s, and probe radius $a=0.243~\mu$m.

% ======================================================================
\section{Finite-size-scaling protocol and results}
% ======================================================================
We repeat the analysis at $L\in\{50,100,200,300,500,750\}$ for
$\kap\in\{0,0.6,0.8,0.97,1.0\}$. Because the finite-size crossover time grows as $L^2$, the steps
per walker are scaled as $L_W(L)=10^6(L/500)^2$, keeping the fit window within the anomalous
regime at every size. From the size-dependent gel points $\pcp(L)$ (Table~\ref{tab:pcL}) and
logistic transition widths $w(L)$ we estimate the correlation-length exponent three ways:
(i) the threshold shift $\pcp(L)=\pcp(\infty)+aL^{-1/\nu}$; (ii) the width scaling
$w(L)\sim L^{-1/\nu}$; and (iii) the data collapse of $\alpha(p,L)$ against
$(p-\pcp(\infty))L^{1/\nu}$, scored by the residual to a master curve in $\alpha$-space (bounded,
so axis rescaling cannot spuriously improve the score). The two smallest sizes ($L=50,100$)
carry the strongest corrections to scaling and are excluded from all exponent fits; uncertainties
are by jackknife over the fitted sizes.

Because $\pcp(L)$ is already stationary for $L\ge200$ (Table~\ref{tab:pcL}), the shift fit
determines $\pcp(\infty)$ robustly but has little leverage on $\nu$; the exponent is therefore
taken from the width and collapse estimators, which do have leverage. These agree on
\begin{equation}
  \nueff\approx1.6, \qquad \text{independent of } \kap \text{ for } \kap\le0.8,
\end{equation}
as summarized in Table~\ref{tab:nu}. The value differs from the static 3D site-percolation
exponent $\nu=0.88$ because $\pcp$ is defined dynamically (through the probe-diffusion exponent)
rather than through a static order parameter. The strongly clustered $\kap=0.97$ condition sits
near the Eden crossover, where the transition broadens and both estimators become noisy
($\nu\approx2.5$--$2.9$); we do not draw exponent conclusions from it. The Eden limit $\kap=1$
has $\pcp\to1$ and no finite-$L$ crossing within the accessible window, and is treated as a bound.

\begin{table}[h]
\centering
\caption{Size-dependent gel point $\pcp(L)$ and its extrapolation $\pcp(\infty)$
(shift fit, $L\ge200$).}
\label{tab:pcL}
\begin{tabular}{lcccc}
\toprule
$L$ & $\kap=0$ & $\kap=0.6$ & $\kap=0.8$ & $\kap=0.97$ \\
\midrule
50   & 0.734 & 0.754 & 0.773 & 0.881 \\
100  & 0.693 & 0.720 & 0.740 & 0.840 \\
200  & 0.684 & 0.706 & 0.725 & 0.815 \\
300  & 0.683 & 0.705 & 0.724 & 0.812 \\
500  & 0.681 & 0.702 & 0.721 & 0.805 \\
750  & 0.682 & 0.703 & 0.724 & 0.804 \\
\midrule
$\infty$ & \textbf{0.681} & \textbf{0.702} & \textbf{0.722} & \textbf{$\approx$0.80} \\
\bottomrule
\end{tabular}
\end{table}

\begin{table}[h]
\centering
\caption{Effective correlation-length exponent from data collapse and transition-width scaling
($L\ge200$; jackknife uncertainties). The shift fit does not constrain $\nu$ (see text). The
$\kap=0.97$ row is a noisy crossover, excluded from the reported $\nueff$.}
\label{tab:nu}
\begin{tabular}{lcc}
\toprule
$\kap$ & $\nu$ (collapse) & $\nu$ (width) \\
\midrule
0.0  & 1.62 & $1.77\pm0.21$ \\
0.6  & 1.64 & $1.49\pm0.05$ \\
0.8  & 1.67 & $1.69\pm0.03$ \\
0.97 & 2.50 & $2.94\pm0.47$ \;(crossover) \\
\bottomrule
\end{tabular}
\end{table}

% ======================================================================
\section{Fixed-$L$ gel-point curve $\pcp(\kap)$}
% ======================================================================
At $L=500$, an 18-value $\kap$-grid resolves the full gel-point curve (Table~\ref{tab:kappa};
mean over $N_s=3$ replicates, $\alpha=0.5$ crossing). The curve is monotonic and consistent with
the thermodynamic-limit anchors of Table~\ref{tab:pcL}. Two regimes are separated near a $\sim1\%$
seed fraction ($\kap\approx0.99$): a gradual rise for $\kap\lesssim0.99$ (many nuclei,
quasi-random morphology) and a steep rise toward unity for $\kap\gtrsim0.99$ (single-cluster
Eden morphology).

\begin{table}[h]
\centering
\caption{Gel point $\pcp$ at $L=500$ for all 18 nucleation densities.}
\label{tab:kappa}
\begin{tabular}{cc@{\hspace{2em}}cc}
\toprule
$\kap$ & $\pcp(L{=}500)$ & $\kap$ & $\pcp(L{=}500)$ \\
\midrule
0.00 & 0.683 & 0.94  & 0.773 \\
0.20 & 0.687 & 0.97  & 0.806 \\
0.40 & 0.692 & 0.98  & 0.824 \\
0.60 & 0.703 & 0.99  & 0.852 \\
0.70 & 0.709 & 0.995 & 0.880 \\
0.80 & 0.721 & 0.996 & 0.888 \\
0.85 & 0.733 & 0.997 & 0.902 \\
0.88 & 0.744 & 0.998 & 0.910 \\
0.91 & 0.751 & 1.00  & $\to1$ (Eden bound) \\
\bottomrule
\end{tabular}
\end{table}

% ======================================================================
\section{Frequency-domain (GSER) spectra}
% ======================================================================
MSD curves near $\pcp$ are converted to $G'(\omega)$, $G''(\omega)$, and $\delta(\omega)$ via the
GSER. Because the MSD is close to a single power law over the anomalous window, the spectra are
power laws with frequency-independent $\delta=90\alpha$ over the accessible range
$\omega\approx0.4$--$4~\mathrm{rad\,s^{-1}}$. At the gel point $G'\approx G''$ and
$\delta\approx45^{\circ}$ for every $\kap$ (Fig.~\ref{fig:gser}); below it $\delta>45^{\circ}$
($G''>G'$, viscous), above it $\delta<45^{\circ}$ ($G'>G''$, elastic). The current spectra are
single-replicate per condition; replication would tighten the modulus estimates.

% ======================================================================
\section{Robustness: alternative growth rules}
% ======================================================================
Connectivity-preserving growth from many independently seeded nuclei per step (an
adjacency-templated rule) also elevates the gel point; for the specific rule studied,
$\pcp(\infty)\approx0.93$. Six-neighbor and twenty-six-neighbor adjacency give indistinguishable
gel points, showing the shift is driven by the connectivity constraint itself, not by the precise
adjacency rule. This confirms that morphology control of the gel point is generic to
connectivity-preserving growth and not an artifact of the $\kap$ construction; the distinct
$\pcp$ of this rule simply reflects a different point in growth-morphology space.

% ======================================================================
\section{Computational implementation and availability}
% ======================================================================
Simulations were implemented in MATLAB with \texttt{parfor} parallelization over walkers.
Lattices are generated cumulatively as described above; the random-walk engine tracks both wrapped
coordinates (for neighbor lookup) and unwrapped coordinates (for the MSD), and the ensemble MSD
is accumulated online by parallel reduction. The finite-size-scaling driver runs each $\kap$
condition at all six sizes with $L_W=10^6(L/500)^2$, $N_W=3000$, $N_s=3$, and the analysis
extracts $\pcp(L)$, $w(L)$, and the three exponent estimates with jackknife uncertainties, fits
restricted to $L\ge200$. Code and processed data are available from the authors on request
[or: at \texttt{[repository link]}].

% ---------------- SM figures ----------------
\begin{figure}[h]
  \centering
  \includegraphics[width=0.7\linewidth]{fig0_schematic_growth.png}
  \caption{Schematic of growth morphologies at fixed occupation, from random ($\kap=0$) through
    clustered intermediate $\kap$ to single-seed (Eden, $\kap=1$) growth, illustrating how compact
    growth leaves a more connected void and thus a higher gel point.}
  \label{fig:schematic}
\end{figure}

\begin{figure}[h]
  \centering
  \includegraphics[width=0.9\linewidth]{fig3_gser_spectra.png}
  \caption{GSER-derived $G'(\omega)$, $G''(\omega)$ and phase angle $\delta(\omega)$ for
    occupation densities straddling $\pcp$. At the gel point $G'\approx G''$ and
    $\delta\approx45^{\circ}$ (the Winter--Chambon signature).}
  \label{fig:gser}
\end{figure}

\bibliographystyle{apsrev4-2}
\bibliography{references_enhanced}

\end{document}
